\documentclass[11pt]{article}

\usepackage[margin=1in]{geometry}
\usepackage{amsmath,amssymb}
\usepackage{booktabs}
\usepackage{graphicx}
\usepackage{xcolor}
\usepackage{listings}
\usepackage[round]{natbib}
\usepackage{url}
\usepackage[hidelinks]{hyperref}
\usepackage{textcomp}

\bibliographystyle{plainnat}

\title{A Six-Layer Evaluation Stack and a Portfolio-Aware Kill Criterion\\
       for Crypto-Perp Trading Strategy Research at Laptop Scale}

\author{David Goh\\
\texttt{davidcagoh@gmail.com}}
\date{May 2026}

\begin{document}
\maketitle

\begin{abstract}
We report a methodology --- not a strategy --- for evaluating systematic trading
strategies on crypto perpetual futures when the practitioner is constrained to
small trade counts ($N{<}200$), laptop-scale compute, and a single execution
substrate. We argue that single-metric leaderboards (Calmar, Sharpe, or even
the Deflated Sharpe Ratio of \citet{lopezdeprado2014dsr}) over-reject genuine
strategies and under-reject overfit ones at this scale, and propose a
\emph{six-layer evaluation stack} in which each layer addresses a failure mode
of the one above: (L1) raw return, (L2) risk-adjusted return, (L3)
sample-size awareness, (L4) multiple-testing deflation, (L5) tail and path
shape, and (L6) portfolio diversification through a Marginal Diversification
Benefit (MDB) triplet. The load-bearing contributions are: a
\emph{portfolio-aware} hard-MDD kill criterion (K1) that corrects the
over-rejection induced by single-strategy pre-registration when adding a
candidate to an existing book, anchored on combined-book MDD, robust MDB, and
maximum pairwise Pearson; a demonstration that \emph{continuous position
shrinkage} along a bull-return / bear-drawdown Pareto frontier dominates
binary regime gates, exhibited via a V1/V2/V3 sweep of an HMM-and-slope
strategy; and a candid treatment of DSR over-rejection at $N{<}200$ with
kurtosis 5--680 typical of trend strategies with fat right tails. We apply
the framework to a candidate book of two strategies (T3 \texttt{SmaRegime180},
R$\wedge$T2 \texttt{HmmSmaSlopeV2}) with pairwise Pearson correlation 0.07,
combined-book risk-parity MDD 2.12\%, and robust MDB +0.55. We name the
ablations: pairs/cointegration (X1) and funding mean-reversion (F1) fall out
of the framework cleanly. The paper is paired with a forward held-out window
(Binance 2026-06 $\to$ 2026-12) that has not been downloaded and will not be
until after submission; results on it will be reported verbatim.
\end{abstract}

\section{Introduction}
\label{sec:intro}

Backtest-driven strategy research at laptop scale is a constrained
optimisation problem: thin trade counts, a single execution substrate (here,
Binance USDT-margined perpetual futures via Freqtrade~\citep{freqtrade}), a
multi-year but still finite history, and the practitioner's own search
process generating selection bias. The pre-registration discipline of
\citet{bailey2014pbo} mitigates the last of these; the Deflated Sharpe Ratio
of \citet{lopezdeprado2014dsr} addresses multiple-testing inflation. Both are
necessary; neither, applied alone, is sufficient. In our setting (5.5 years
of Binance perpetuals, 5 coins, $N$ closed trades typically in $[30, 700]$,
returns with sample skew up to $+24$ and excess kurtosis up to $+680$), DSR
flags every candidate strategy as ``noise.'' If we gate on DSR alone, the
project is over. If we ignore it, we replay the same overfitting failures
that motivated its introduction.

This paper documents what we did instead. We construct a six-layer stack
(\S\ref{sec:stack}) in which each layer addresses a failure mode of the one
above. Critically, we treat L4 (DSR) as a \emph{humility check} rather than a
binding gate, and add L5 (tail and path shape) and L6 (portfolio
diversification) as load-bearing gates. We then formalise a
\emph{portfolio-aware} hard-MDD kill criterion (\S\ref{sec:portfolio-k1}) ---
the single most important methodological correction in the paper --- that
prevents a pre-registered single-strategy kill threshold from over-rejecting a
candidate whose drawdown is uncorrelated with the existing book.

\subsection{Contributions}
\label{sec:contributions}

\begin{enumerate}
  \item \textbf{A six-layer evaluation stack} (L1--L6) in which each layer
        addresses a specific failure mode of the previous layer. The stack is
        portable: the layer count is fixed, but the gate thresholds at L2,
        L3, L5, and L6 are calibrated to the substrate.
  \item \textbf{A portfolio-aware K1 kill rule} (\S\ref{sec:portfolio-k1})
        with four clauses --- combined-book MDD, robust MDB triplet
        (equal-weight, risk-parity, mean-var), max pairwise Pearson, and a
        $2\times K1$ hard cap. We show that the standalone K1 of
        \citet{decision_004} over-rejects when applied as a portfolio
        admission rule, and the corrected gate admits the second strategy in
        our candidate book on principled grounds.
  \item \textbf{Continuous shrinkage along a Pareto frontier} as a dominant
        alternative to binary regime gates. A V1/V2/V3 sweep of the
        HMM-and-slope strategy (\S\ref{sec:shrinkage}) traces a continuous
        bull-return / bear-drawdown frontier rather than collapsing to a
        single operating point.
  \item \textbf{Empirical scaffolding} for a candidate two-strategy book
        \{T3 \texttt{SmaRegime180}, R$\wedge$T2 \texttt{HmmSmaSlopeV2}\}
        on Binance 5-coin perpetuals 2020-09 $\to$ 2026-05, with pairwise
        Pearson 0.07 and robust MDB-rp $+0.55$. Five strategy families that
        \emph{fail} the stack (X1 pairs, F1 funding MR, C1 funding carry, R2
        HmmRegime4Rolling, R$\wedge$C1 HmmCarry conjunction) function as
        ablations of the framework.
\end{enumerate}

\subsection{What this paper is not (and why)}

This paper is not a strategy proposal. The candidate book is the
\emph{worked example} of the methodology, not its endpoint. It is also not
yet an out-of-sample claim: the forward window 2026-06 $\to$ 2026-12 has
intentionally not been downloaded, and the 30-day live paper-trade dry-run is
unfinished. We are explicit about both gaps (\S\ref{sec:threats}). Finally,
this is not a critique of \citet{lopezdeprado2014dsr}: DSR is the correct
gate at institutional $N$; it is the wrong gate at laptop $N$ with heavy
right-tailed distributions. The contribution is a framework that names the
regime in which DSR is informative versus binding.

\paragraph{Reproducibility.} All backtests use Freqtrade with
\texttt{--fee 0.00035} (Hyperliquid taker, applied uniformly across the
universe for cross-strategy comparability), public Binance perpetuals OHLCV
and funding data via \texttt{ccxt}, and the common evaluation substrate
documented in \citet{decision_005}. Source code, decision documents, and
per-strategy result cards are in the companion repository \texttt{algo-traders}.

\section{Related Work}
\label{sec:related}

\subsection{Backtest evaluation and overfitting}

The line of work descending from \citet{bailey2014pbo} and
\citet{lopezdeprado2014dsr} formalises two distinct failure modes of
naive backtest evaluation. The Probability of Backtest Overfitting (PBO)
captures the selection-bias problem: an analyst running $K$ trials and
reporting the best Sharpe systematically overstates out-of-sample
performance. The Deflated Sharpe Ratio (DSR) corrects the observed Sharpe
for: (a) the number of trials, (b) the cross-trial Sharpe variance, and (c)
the skew and kurtosis of the strategy returns. The skew/kurtosis correction
is the term that bites at our scale --- with kurtosis 5--680 our raw Sharpe
estimates are aggressively shrunken even before the multiple-testing
adjustment, and the resulting DSR is below conventional significance for
every standalone strategy we tested. We retain DSR as Layer 4 (\S\ref{sec:l4})
but do not gate admission on it.

\citet{simhi2026oldhabits} is the related warning from the LLM-assisted
strategy generation regime: generators trained on retail trading content
recycle discredited indicators (e.g.\ RSI thresholds, MACD crossovers)
because they dominate the training corpus. The lesson translates: a layered
gate with explicit second- and third-moment terms is robust to the
``looks-like-a-strategy'' failure mode in a way that a single-metric gate
is not.

\subsection{Crypto-perp microstructure and funding}

\citet{le2026fundingmm} provides the funding-rate Ornstein-Uhlenbeck
calibration we adopt as a default prior for funding-driven strategies, with
empirical half-life $\approx 8$h on Hyperliquid majors.
\citet{badawi2026funding4h} and \citet{inan2025fundingpred} establish
short-horizon predictability of funding rates that justifies a 4h cadence
funding-aware strategy; \citet{twotiered2026} and \citet{dexcarry2025} document
the CEX/DEX two-tier structure that determines whether funding is a return
source or a cost. \citet{sar2026} provides Hyperliquid-specific slippage
distributions that inform our fee assumption (\texttt{--fee 0.00035} is the
Hyperliquid taker; applying it uniformly to Binance backtests is a
conservative bias). The funding mean-reversion strategy F1
(\S\ref{sec:ablations}) is the failed test of the funding-extreme thesis on
our substrate; it falls out of the framework cleanly.

\subsection{Regime detection}

\citet{pang2026hmmbtc} applies Gaussian HMMs to Bitcoin daily returns;
\citet{wassersteinhmm2026} extends to a Wasserstein-distance-based variant
that is more robust to non-Gaussian state emissions. Our R-family strategies
(\texttt{HmmRegime4}, \texttt{HmmRegime4Rolling}) implement a standard
4-state Gaussian HMM on 1h returns. \citet{backtestlive2026} documents the
``regime-timing'' gap between backtest and live performance, which motivates
treating $\sim$ (look-ahead) HMM variants as unbeatable ceilings, not
strategies. \citet{trendsreversion2025} establishes the 1h--4h timescale as
the boundary at which trend and mean-reversion structures co-exist; this is
the scale at which we run.

\section{The Six-Layer Evaluation Stack}
\label{sec:stack}

We define six layers, each posed as a question and answered by an explicit
metric. Each layer addresses a specific failure mode of the previous one
(Table~\ref{tab:layers}).

\begin{table}[h]
\centering
\small
\begin{tabular}{cllp{5cm}}
\toprule
Layer & Question & Headline metric & Fixes failure of \\
\midrule
L1 & Profitable? & CAGR & --- \\
L2 & Risk-adjusted? & Calmar, Sharpe & L1: scale-free \\
L3 & Sample-size aware? & SQN & L2: thin samples \\
L4 & Multiple-testing? & DSR & L3: trial inflation \\
L5 & Tail/path? & Ulcer, Martin, kurtosis, CVaR$_{5\%}$ & L2: Gaussian lies \\
L6 & Portfolio-additive? & MDB-rp, corr & L2--L5: redundancy \\
\bottomrule
\end{tabular}
\caption{The six-layer evaluation stack. Each layer fixes a failure mode of
the previous layer; the stack as a whole gates admission to the candidate
book.}
\label{tab:layers}
\end{table}

The reading rule we follow: a strategy is admitted to paper-trade if it is
positive on L2--L3--L5 \emph{and} has robust positive MDB-rp at meaningful
magnitude against the current book. L4 (DSR) is reported but not gating;
\S\ref{sec:l4} explains why.

\subsection{Layer 1: Return}
L1 is total return (or CAGR). It is the layer below which nothing matters and
above which nothing is decided. We report it but never gate on it.

\subsection{Layer 2: Risk-adjusted return}
\label{sec:l2}
\textbf{Calmar} (CAGR over maximum drawdown) and \textbf{Sharpe} (annualised
return over annualised volatility, 365-day basis). Calmar is our headline
because drawdown rather than volatility is the binding constraint at laptop
scale --- a 20\% drawdown ends the experiment regardless of long-run
expectation. Sharpe is reported for cross-paper comparability.

\subsection{Layer 3: Sample-size awareness}
\label{sec:l3}
The System Quality Number,
\begin{equation}
\mathrm{SQN} \;=\; \frac{\bar{r}}{\sigma_r} \cdot \sqrt{N},
\end{equation}
where $\bar{r}$ and $\sigma_r$ are the mean and standard deviation of
\emph{per-trade} returns and $N$ is the trade count. SQN penalises strategies
whose Calmar is driven by a small number of trades. In our universe,
$\mathrm{SQN} \geq 1.7$ on $N \geq 80$ trades is the empirical threshold below
which standalone Calmar should not be trusted as a point estimate.

\subsection{Layer 4: Multiple-testing deflation --- and when it over-rejects}
\label{sec:l4}
The Deflated Sharpe Ratio of \citet{lopezdeprado2014dsr} is
\begin{equation}
\mathrm{DSR} \;=\; \Phi\!\left(
  \frac{(\widehat{\mathrm{SR}} - \mathrm{SR}_0)\sqrt{N-1}}
       {\sqrt{1 - \hat{\gamma}_3 \widehat{\mathrm{SR}} + \frac{\hat{\gamma}_4 - 1}{4}\widehat{\mathrm{SR}}^2}}
\right),
\label{eq:dsr}
\end{equation}
where $\widehat{\mathrm{SR}}$ is the observed Sharpe, $\mathrm{SR}_0$ is a
benchmark adjusted for the number of trials and cross-trial Sharpe variance,
$\hat{\gamma}_3$ and $\hat{\gamma}_4$ are sample skew and kurtosis of
returns, and $\Phi$ is the standard normal CDF.

Two terms in the denominator dominate at our scale. First, $\hat{\gamma}_4$
ranges from 5 to 680 across our strategies; the kurtosis correction shrinks
the test statistic by a factor of 3--20. Second, $\sqrt{N-1}$ with
$N \in [30, 700]$ does not overcome the variance inflation. The empirical
result, documented in our DSR sweep~\citep{result_dsr_analysis}, is that
\emph{every} standalone strategy receives DSR below conventional
significance, including strategies with Calmar $> 8$ on $> 90$ trades across
two bull/bear cycles. Gating on DSR alone retires the entire project.

We adopt the following rule: DSR is computed and reported on every result
card, but admission is gated on L2, L3, L5, and L6 jointly. DSR functions as
a humility check --- a signal to look harder at L5 (tail shape) and L6
(portfolio additivity), not a binary kill. This is the empirical analogue of
the survival-of-the-quant comment in \citet{davies2025survival}: a
practitioner who only ships strategies with DSR $> 0.5$ at $N=100$ with
kurtosis $50$ ships nothing.

\subsection{Layer 5: Tail and path shape}
\label{sec:l5}

L2 reports the worst single drawdown; L5 reports the \emph{shape} of the
underwater experience.

The \textbf{Ulcer Index} of the equity curve $\{W_t\}$ is
\begin{equation}
\mathrm{Ulcer} \;=\; \sqrt{\frac{1}{T} \sum_{t=1}^{T}
\left(\frac{W_t - \max_{s \leq t} W_s}{\max_{s \leq t} W_s}\right)^2 \cdot 100^2}.
\end{equation}
Ulcer is path-aware: a strategy with a single deep drawdown that recovers
quickly has low Ulcer; a strategy with a chronic shallow drawdown has high
Ulcer despite a smaller MDD. The Martin ratio is
$\mathrm{Martin} = \mathrm{CAGR} / \mathrm{Ulcer}$.

We additionally report sample \textbf{skew} ($\hat{\gamma}_3$), excess
\textbf{kurtosis} ($\hat{\gamma}_4$), tail ratio
$|P_{95}|/|P_5|$, and CVaR$_{5\%} = \mathbb{E}[r \mid r \leq P_5]$ of daily
log-returns. Slope-gated strategies in our universe have $\hat{\gamma}_3 \in
[+9, +21]$ and $\hat{\gamma}_4 \in [+110, +470]$. Sharpe overstates these by
a Jensen-inequality margin that L5 makes visible.

\paragraph{Worked comparison.} T3 (\texttt{SmaRegime180}) has MDD 2.21\%,
Ulcer 1.30, Pain 1.12 --- a tight, fast-recovering path. R2
(\texttt{HmmRegime4Rolling}) has MDD 7.65\% with Ulcer 4.01, Pain 3.47 ---
chronic underwater. The MDD gap (2.21\% vs 7.65\%) understates the path-shape
gap; L5 makes the gap binding.

\subsection{Layer 6: Portfolio diversification}
\label{sec:l6}

L1--L5 evaluate strategies in isolation. L6 evaluates them as additions to a
current book.

Let $B = \{s_1, \ldots, s_k\}$ be the current book and $c$ a candidate.
Define the equity curve of book $X$ under weighting scheme $w$ as
$E_w(X)$ and the Sharpe as $\mathrm{SR}(\cdot)$. The
\textbf{Marginal Diversification Benefit} is
\begin{equation}
\mathrm{MDB}_w(c, B) \;=\; \mathrm{SR}(E_w(B \cup \{c\})) \;-\; \mathrm{SR}(E_w(B)).
\end{equation}

We compute MDB under three weighting schemes:
\begin{itemize}
  \item \textbf{MDB-eq:} equal-weight, $w_i = 1/|B \cup \{c\}|$.
  \item \textbf{MDB-rp:} risk-parity. Trailing 90-day return volatility
        $\sigma_i$ per strategy; weights $w_i \propto 1/\sigma_i$, renormalised.
        This is the leaderboard headline.
  \item \textbf{MDB-mv:} mean-variance, long-only quadratic optimisation.
        Reported as an upper bound; known to be unstable at small $N$.
\end{itemize}

A candidate is \textbf{robustly} diversifying iff MDB $> 0$ under all three
schemes. The robustness flag is the load-bearing protection against
weighting-artifact false positives: a strategy whose only positive MDB comes
from the mean-variance scheme is being rescued by optimiser noise, not by
genuine diversification.

\paragraph{Headline result.} The candidate book \{T3, R$\wedge$T2\} has
pairwise Pearson correlation 0.07; R$\wedge$T2's MDB-rp against
\{T3\} alone is $+0.55$, robust positive under all three schemes.

\section{Portfolio-Aware Kill Criteria}
\label{sec:portfolio-k1}

\subsection{The single-strategy hard kill (K1)}
\label{sec:k1-standalone}

Decision document 004~\citep{decision_004} pre-registers a hard-MDD kill
threshold for the T3 \texttt{SmaRegime180} strategy:
\begin{equation}
K1_{\mathrm{standalone}} \;=\; 5.5\% \;\approx\; 3 \cdot \mathrm{IS\text{-}MDD}(T3, \text{Hyperliquid bear}).
\end{equation}
The $3\times$ multiplier captures the empirical observation that live
maximum drawdown typically exceeds in-sample MDD by a factor of three under
regime shift. This is the standalone gate: if any single strategy in the
book breaches 5.5\% live MDD, it is retired.

\subsection{Why standalone K1 over-rejects portfolio additions}

R$\wedge$T2 \texttt{HmmSmaSlopeV2} has standalone MDD 6.05\% on the common
window --- 0.55 percentage points above K1. Under a strict reading,
R$\wedge$T2 is killed and the book collapses to a single strategy. Two
implicit assumptions in the standalone K1 derivation do not transfer:

\begin{enumerate}
  \item \textbf{Single-window scope.} The 1.74\% IS anchor was measured on
        one 6-month BTC-only bear. On the 5.5-year 5-coin Binance common
        window T3's \emph{own} MDD widens to 2.21\% --- the same strategy on
        a more demanding substrate. Consistent application of the $3\times$
        rule places K1 at $\approx 6.6\%$ in the common-window basis.
  \item \textbf{Single-strategy scope.} A 6.05\% drawdown that occurs at a
        \emph{different time} from T3's drawdowns is a fundamentally
        different risk object from a 6.05\% drawdown that compounds T3's. At
        pairwise Pearson 0.07, the combined-book MDD under risk-parity
        weighting is bounded well below the sum.
\end{enumerate}

\subsection{The portfolio-aware gate}
\label{sec:gate}

A candidate strategy $c$ may enter a book $B$ (where $B$ already contains at
least one standalone-K1-compliant ``anchor'' strategy) if \emph{either} the
standalone path or the portfolio path is satisfied:
\begin{align}
\textbf{Standalone path:} \quad &\mathrm{MDD}(c) \leq K1_{\mathrm{standalone}} \\
\textbf{Portfolio path:} \quad &\mathrm{MDD}(c) \leq K1_{\mathrm{hard\,cap}}
                                  = 2 \cdot K1_{\mathrm{standalone}} \nonumber \\
                          \wedge\;\; &\mathrm{MDD}(B \cup \{c\}) \leq K1_{\mathrm{book}}
                                  = K1_{\mathrm{standalone}} \nonumber \\
                          \wedge\;\; &\mathrm{MDB}_w(c, B) > 0 \;\;\forall w \in \{\mathrm{eq}, \mathrm{rp}, \mathrm{mv}\} \nonumber \\
                          \wedge\;\; &\mathrm{MDB}_{\mathrm{rp}}(c, B) \geq 0.30 \nonumber \\
                          \wedge\;\; &\max_{s \in B} |\rho(c, s)| < 0.85.
\label{eq:gate}
\end{align}

Each clause has a specific failure mode it addresses. The hard cap
$2 \cdot K1$ prevents arbitrarily large standalone drawdowns from being
``rescued'' by trivially low risk-parity weight (\S\ref{sec:bounds}). The
book-MDD clause ensures the portfolio's own drawdown still respects the
original ceiling. The robust MDB triplet prevents weighting-artifact
admission. The MDB-rp magnitude threshold $\geq 0.30$ requires a meaningful
diversification benefit, not noise-level positivity. The Pearson clause
prevents statistical duplicates (e.g.\ R$\wedge$T1/V2/V3 collectively, with
mutual correlations $0.96\text{--}1.00$, cannot all be in the book).

\subsection{R\texorpdfstring{$\wedge$}{ }T2 walkthrough}

Applied to candidate R$\wedge$T2 against book $B = \{T3\}$:

\begin{table}[h]
\centering
\small
\begin{tabular}{lrrc}
\toprule
Clause & Required & Observed & Pass \\
\midrule
$\mathrm{MDD}(c) \leq 5.5\%$ & 5.5\% & 6.05\% & \texttimes \\
$\to$ falls through to portfolio path & & & \\
$\mathrm{MDD}(c) \leq 11.0\%$ & 11.0\% & 6.05\% & \checkmark \\
$\mathrm{MDD}_{\mathrm{rp}}(B \cup \{c\}) \leq 5.5\%$ & 5.5\% & 2.12\% & \checkmark \\
robust MDB $> 0$ (eq, rp, mv) & all $> 0$ & all $> 0$ & \checkmark \\
$\mathrm{MDB}_{\mathrm{rp}} \geq 0.30$ & 0.30 & $+0.55$ & \checkmark \\
$\max \rho(c, B) < 0.85$ & 0.85 & 0.07 & \checkmark \\
\bottomrule
\end{tabular}
\caption{Portfolio-aware gate, applied to R$\wedge$T2 entering book \{T3\}.}
\label{tab:gate-applied}
\end{table}

The combined-book MDD of 2.12\% is computed under risk-parity weights at
window endpoint $\{T3: 0.685, \, R\wedge T2: 0.315\}$ on the 2054-day common
window. The two strategies' drawdowns occur at different times, and the
risk-parity weighting caps R$\wedge$T2's contribution; the combined book sits
at less than 40\% of the original gate.

\subsection{Bounds on the relaxation}
\label{sec:bounds}

The portfolio-aware gate is not a license for arbitrarily larger standalone
drawdowns. Three boundary conditions are explicit:

\begin{enumerate}
  \item \textbf{Hard cap $K1_{\mathrm{hard\,cap}} = 2 \cdot K1_{\mathrm{standalone}}$.}
        Above this, no portfolio justification rescues the candidate. At
        $\mathrm{MDD} > 11\%$, the only way risk-parity can neutralise the
        strategy's contribution is to drive its weight toward zero --- a
        ``phantom entry'' that inflates the leaderboard without changing
        live P\&L. R2x (\texttt{HmmRegime4Rolling} 5-coin, MDD 21.47\%) is
        the worked example: technically MDB-positive but excluded for
        precisely this reason.
  \item \textbf{No daisy-chaining.} A strategy admitted via the portfolio
        path does not itself become an anchor. The book must always contain
        at least one standalone-K1-compliant strategy.
  \item \textbf{Cross-cycle requirement.} Portfolio-path admission requires
        the common window to span at least one bull and one bear regime
        (here, two of each). A K1 breach measured on a single regime cannot
        be rescued by an MDB that is itself regime-conditional.
\end{enumerate}

\section{Empirical Application}
\label{sec:empirical}

\subsection{Data}
Binance USDT-margined perpetual futures, 5 coins (BTC, ETH, SOL, AVAX, DOGE),
4h and 1h bars, common window 2020-09-23 $\to$ 2026-05-09 (2054 days, 2 bulls
and 2 bears). Fees \texttt{--fee 0.00035} (Hyperliquid taker, applied
uniformly). Funding rates downloaded for the full 5.5-year window in May 2026
(superseding an earlier 2.3-year truncated parquet). Data sourcing via
\texttt{ccxt}; pipeline in the companion repository.

\subsection{Strategy universe}

Seven families, distinguished by signal source and conjunction structure:
\textbf{T} (trend, SMA + slope), \textbf{R} (regime, HMM posterior),
\textbf{C} (carry, funding rate), \textbf{R$\wedge$T} (regime $\times$ trend
conjunction), \textbf{R$\wedge$C} (regime $\times$ carry), \textbf{X}
(cross-sectional / pairs), \textbf{F} (funding mean-reversion). Each
strategy carries a pre-registered kill-criteria document (\S\ref{sec:bounds},
\citep{decision_004, decision_005}).

\subsection{The candidate book \{T3, R$\wedge$T2\}: per-layer evidence}

% Common-window leaderboard, Binance perp 2020-09-23 -> 2026-05-09 (5.5y).
% Risk-parity MDB computed against the candidate book {T3, R^T2}.
\begin{table*}[t]
\centering
\small
\begin{tabular}{llrrrrrrrl}
\toprule
Code & Strategy & Calmar & SQN & MDD\% & Ulcer & Martin & corr$\to$T3 & MDB-rp & Status \\
\midrule
T3      & SmaRegime180          & 8.76  & 1.78 & 2.21  & 1.30  & +2.51 & (book) & (book) & \textbf{$\star$} \\
R$\wedge$T2 & HmmSmaSlopeV2     & 30.23 & 2.73 & 6.05  & 2.87  & +7.41 & (book) & (book) & \textbf{$\star$} \\
R$\wedge$T3 & HmmSmaSlopeV3     & 27.28 & 2.79 & 6.91  & 3.30  & +6.58 & +0.07  & $-0.023$ & $\blacktriangle$ \\
R$\wedge$T1 & HmmSmaSlope (V1)  & 25.01 & 2.97 & 8.21  & 3.82  & +6.00 & +0.07  & $+0.012$ & $\blacktriangle$ \\
R1$\sim$ & HmmRegime4 (look-ahead) & 9.16 & 3.88 & 2.94 & 1.09 & +4.24 & 0.00 & $+0.07$ & $\sim$ \\
T2      & SmaRegime720          & 5.39  & 1.74 & 3.57  & 1.75  & +1.82 & +0.65  & $-0.023$ & $\blacktriangle$ \\
T1      & TrendFilter200        & 1.69  & 1.10 & 8.54  & 3.79  & +0.67 & 0.00 & $+0.040$ & $\blacktriangle$ \\
X2      & CrossSectionalMom.    & ---   & ---  & 13.04 & ---   & ---   & 0.00 & $+0.048$ & $\blacktriangle$ \\
R2x     & HmmRegime4Rolling 5c  & 3.79  & 1.66 & 21.47 & 10.25 & +1.15 & $-0.01$ & $+0.069$ & \texttimes \\
R2      & HmmRegime4Rolling BTC & 0.47  & 0.39 & 7.65  & 4.01  & +0.17 & +0.01  & $+0.010$ & \texttimes \\
X1      & PairsZScore (1-leg)   & 2.08  & 0.95 & 4.03  & 0.48  & +0.55 & 0.00 & $-0.898$ & \texttimes \\
X1v2    & PairsZScoreV2 (2-leg) & 0.97  & 0.42 & 2.91  & 0.82  & +0.65 & n/c    & n/c     & \texttimes \\
C1      & FundingCarry 5c (full)& 1.93  & 1.22 & 15.65 & 5.87  & +0.81 & n/c    & n/c     & \texttimes \\
R$\wedge$C1 & HmmCarry 5c       & 0.08  & 0.14 & 35.46 & 9.57  & +0.04 & 0.00 & $-0.000$ & \texttimes \\
R$\wedge$C1$^r$ & HmmCarryReverse 5c & 0.78 & 0.56 & 15.09 & 7.50 & +0.28 & n/c & n/c & \texttimes \\
F1      & FundingExtremeMR 5c   & $-0.32$& $-0.64$& 33.51 & 15.98 & $-0.14$ & n/c & n/c & \texttimes \\
T0      & LongOnlyStrategy      & $-0.37$& $-0.58$& 10.17 & 6.26  & $-0.11$& 0.00 & $-0.009$ & $\cdot$ \\
\bottomrule
\end{tabular}
\caption{Common-window leaderboard. Binance perp, 5-coin universe (BTC/ETH/SOL/AVAX/DOGE), 4h or 1h bars, fees $\texttt{--fee 0.00035}$ uniformly. Status: $\star$ admitted to book; $\blacktriangle$ research frontier; $\sim$ look-ahead upper bound; \texttimes\ killed; $\cdot$ baseline. ``n/c'' = not computed (book composition unchanged).}
\label{tab:leaderboard}
\end{table*}


Table~\ref{tab:layer-evidence} reports the per-layer evidence for the two
admitted strategies.

\begin{table}[h]
\centering
\small
\begin{tabular}{lrr}
\toprule
Layer / Metric & T3 \texttt{SmaRegime180} & R$\wedge$T2 \texttt{HmmSmaSlopeV2} \\
\midrule
Universe & BTC 4h & 5-coin 4h \\
L1 CAGR & $+3.42\%$ & $+21.31\%$ \\
L2 Calmar & 8.76 & 30.23 \\
L2 MDD & 2.21\% & 6.05\% \\
L3 SQN & 1.78 & 2.73 \\
L3 trades & 85 & 616 \\
L5 Ulcer & 1.30 & 2.87 \\
L5 Martin & $+2.51$ & $+7.41$ \\
L5 skew & $+20.4$ & $+24.0$ \\
L5 kurtosis (excess) & $+455$ & $+680$ \\
L6 $\rho \to$ other & 0.07 & 0.07 \\
L6 MDB-rp vs other & $+0.55$ (book) & $+0.55$ (book) \\
\bottomrule
\end{tabular}
\caption{Per-layer evidence for the candidate book. T3 is the
conservative anchor; R$\wedge$T2 is the high-return / low-correlation
complement admitted via the portfolio-aware gate.}
\label{tab:layer-evidence}
\end{table}

\subsection{Ablations: families that fall out of the framework}
\label{sec:ablations}

Five strategy families fail the stack at distinct layers. Each functions as
an ablation: the layer at which they fail identifies the failure mode the
layer was constructed to catch.

\paragraph{X1 (pairs/cointegration).} Single-leg \texttt{PairsZScore}
SOL-DOGE has standalone Calmar 2.08 and MDD 4.03\%. It passes L1--L5 and
the standalone K1. It fails L6 catastrophically:
$\mathrm{MDB}_{\mathrm{rp}} = -0.898$. The low-vol return profile receives
oversized risk-parity weight without enough return to justify it; the book's
Sharpe drops when X1 is added. The two-leg variant X1v2
\citep{result_pairs_v2} reaches the same kill verdict through a different
failure: Freqtrade cannot atomically pair-stop on spread (per-leg stops fire
on per-instrument prices), so the legs are not co-managed and the strategy
is a directional bet on each leg separately. This is a framework limit, not
a strategy failure (\S\ref{sec:framework-limit}).

\paragraph{F1 (funding mean-reversion).} \texttt{FundingExtremeMR} fails L1
and L2: $-11.60\%$ return, Calmar $-0.32$, MDD 33.51\% on the full 5.5-year
window~\citep{result_funding_mr_fullwin}. Counter-funding does not
mean-revert at 4h cadence on Binance majors; the F1 thesis fails on this
substrate.

\paragraph{C1 (funding carry).} \texttt{FundingCarry} 5-coin shows genuine
standalone edge after the full-window funding parquet: $+32.44\%$ return,
Calmar 1.93, MDD 15.65\%~\citep{result_funding_carry_fullwin}. It fails the
portfolio path on the hard cap clause: MDD 15.65\% $> 11\% = 2 K1$. The L6
evaluation is moot --- C1 cannot be in the book regardless of its MDB.

\paragraph{R2 (HmmRegime4Rolling).} Walk-forward HMM on BTC 1h. Calmar 0.47,
MDD 7.65\%, Ulcer 4.01 --- the L5 worked comparison from \S\ref{sec:l5}.
Fails L2 (Calmar) and the standalone K1. Look-ahead variant R1 has Calmar
9.16 on the same substrate; the 0.47-to-9.16 ratio quantifies the look-ahead
absorption. R2 is a clean ablation of the "honest walk-forward versus
look-ahead" gap that motivates pre-registration.

\paragraph{R$\wedge$C1 (HmmCarry conjunction).} The product of an HMM-bull
signal and a funding-negative signal. Hypothesis: independent signals should
intersect to a tighter filter. Empirically the signals are
\emph{anti-complementary} --- HMM is reactive (lags regime change) while
funding is forward-looking (leads). The intersection picks the worst of both
worlds: late-cycle bull lag with early-cycle bear lead. Calmar 0.08, MDD
35.46\%. The reverse-sign variant R$\wedge$C1$^r$ (HMM-bull AND
funding-positive) is directionally correct but still breaches K1 by $2.7\times$
\citep{result_hmm_carry_reverse}. Both confirm that signal-conjunction
arguments require empirical verification of independence, not assertion.

\section{Continuous Shrinkage along the Pareto Frontier}
\label{sec:shrinkage}

\subsection{The V1/V2/V3 sweep}

We constructed three variants of the HMM-and-slope strategy (R$\wedge$T1,
R$\wedge$T2, R$\wedge$T3) differing only in their position-sizing function
applied to the HMM bull-state posterior $p_t$:

\begin{align*}
\textbf{V1 (binary):}    \quad &h_t = \mathbb{1}[p_t > \tau], \\
\textbf{V2 (linear):}    \quad &h_t = \max(0, \, p_t - \tau), \\
\textbf{V3 (sqrt):}      \quad &h_t = \max(0, \, \sqrt{p_t - \tau}).
\end{align*}

Each variant traces a different point on the bull-return / bear-drawdown
Pareto frontier. The full result triplet (Table~\ref{tab:vsweep}) shows the
frontier directly.

\begin{table}[h]
\centering
\small
\begin{tabular}{lrrrr}
\toprule
Variant & Calmar & SQN & MDD\% & Martin \\
\midrule
V1 (binary) & 25.01 & 2.97 & 8.21 & $+6.00$ \\
V2 (linear) & 30.23 & 2.73 & 6.05 & $+7.41$ \\
V3 (sqrt)   & 27.28 & 2.79 & 6.91 & $+6.58$ \\
\bottomrule
\end{tabular}
\caption{V1/V2/V3 sweep. Pairwise correlations 0.96--1.00 (statistically
one strategy); V2 selected as portfolio representative on best Calmar/MDD
pair.}
\label{tab:vsweep}
\end{table}

The three variants have pairwise Pearson correlations 0.96--1.00: they are
\emph{statistically one strategy} measured at three points on a shrinkage
curve. The L6 framework correctly excludes V1 and V3 from the book once V2
is admitted (Pearson clause; Table~\ref{tab:leaderboard}). They remain on
the research frontier as boundary points.

\subsection{Reading the frontier as a continuous knob}

The binary-versus-continuous-sizing question generalises beyond the V-sweep.
Following \citet{ravagnani2026shrinkage}, we treat
\begin{equation}
h_{\mathrm{robust}} = h_{\mathrm{standard}} \cdot \left(1 - \frac{\delta}{\sigma^2}\right)
\end{equation}
as the canonical continuous-shrinkage form, where $\delta$ is the magnitude
of the regime-correction estimate and $\sigma^2$ is its sampling variance.
Binary gates ($h \in \{0, 1\}$) are the limit of this rule as $\sigma^2 \to
0$; they are statistically wasteful when $\sigma^2$ is finite and the
regime signal is noisy. Decision~004 already encodes this in the
\texttt{calmar\_factor}, \texttt{pf\_factor}, \texttt{regime\_factor}
multipliers~\citep{decision_004}; the V-sweep is the empirical confirmation
that the continuous form sits strictly above the binary form on the Pareto
frontier on this substrate.

\section{Threats to Validity}
\label{sec:threats}

\subsection{Held-out window not yet evaluated}
The forward window Binance 2026-06-01 $\to$ 2026-12-31 has intentionally not
been downloaded as of the submission date. This is the single most important
unrun gate. The pre-commitment is explicit: download once after submission,
run the book and the surviving frontier strategies without parameter tweaks,
report results verbatim. No claim in this paper is robust until this gate is
crossed.

\subsection{Laptop-scale $N$ --- DSR versus evidence tension}
We argue (\S\ref{sec:l4}) that DSR over-rejects at $N < 200$ with kurtosis
$> 50$. This is an empirical claim about our substrate, not a theoretical
critique of the DSR derivation. A reviewer might reasonably reverse the
conclusion: \emph{exactly} these strategies are the ones the DSR construction
is designed to catch, and the framework's L5/L6 admissions are simply
laundered overfitting. We do not have a clean refutation; the held-out
window is the strongest available test, and we have not run it. We name the
tension rather than resolve it.

\subsection{Single-substrate framework calibration}
All gate thresholds (K1 = 5.5\%, hard cap = 11\%, MDB-rp magnitude floor
0.30, Pearson ceiling 0.85) are calibrated on Binance 5-coin perpetuals
2020-09 $\to$ 2026-05. The layer structure is portable; the specific scalars
are substrate-conditional. Application to a different exchange, asset class,
or time horizon requires recalibration via the same pre-registration
discipline that produced the original values.

\subsection{Framework limit: Freqtrade cannot atomically pair-stop on spread}
\label{sec:framework-limit}
The X1v2 two-leg PairsZScore result~\citep{result_pairs_v2} surfaces a
specific limit of the substrate. Freqtrade's per-leg stop-loss machinery
fires on each instrument's own price, not on the joint spread; a divergence
that widens one side's price closes that leg independently of the joint
$z$-score, breaking market-neutrality. This is a framework limit, not a
strategy failure: a true pairs evaluation requires a different execution
substrate (Backtrader, Lean, or custom paired risk management above
Freqtrade). The X1 family verdict in this paper is therefore narrower than
``pairs do not work on crypto majors'': it is ``pairs cannot be
faithfully evaluated within this substrate.''

\subsection{LLM-recycled strategy templates}
Following \citet{simhi2026oldhabits}: this paper's strategies were
hand-coded, but the family taxonomy and metric choices were informed by
LLM-assisted literature triage. The risk that load-bearing methodological
choices reflect training-corpus consensus rather than principled derivation
is non-zero. The L4 humility check and the explicit pre-registration
template are the available mitigations.

\section{Future Work}
\label{sec:future}

\subsection{Forward held-out window 2026-06 $\to$ 2026-12 (pre-registered, frozen)}
Run the candidate book and the surviving research-frontier strategies on
2026-06-01 $\to$ 2026-12-31, one-shot, no parameter tweaks. Report verbatim.
This is the single highest-priority follow-up.

\subsection{30-day live paper-trade dry-run}
Run \{T3, R$\wedge$T2\} on Hyperliquid in paper-trade mode for 30 days
against the decision-004 kill criteria before any real capital. This catches
the backtest-to-live divergence documented in
\citet{backtestlive2026} that the held-out window cannot.

\subsection{Per-coin signed-funding extension}
The R$\wedge$C family (HMM conjunction with funding) is killed in both
sign directions on a single-rule basis. AVAX shows opposite sign to
BTC/ETH/SOL/DOGE in our data, suggesting per-coin signed-funding learning is
the live extension. We leave this to future work.

\section{Conclusion}
\label{sec:conclusion}

We have presented a six-layer evaluation stack for crypto-perp trading
strategy research at laptop scale, a portfolio-aware kill criterion that
corrects single-strategy over-rejection, an empirical confirmation that
continuous shrinkage dominates binary gates along the relevant Pareto
frontier, and a candid treatment of where DSR over-rejects at our $N$ and
kurtosis. The candidate book \{T3, R$\wedge$T2\} with pairwise Pearson 0.07
and combined-book risk-parity MDD 2.12\% is the worked example; the five
ablations (X1, F1, C1, R2, R$\wedge$C1) name the layers at which strategy
families fall out.

The paper is incomplete by design: the forward held-out window and the live
paper-trade dry-run are the two binding tests we have not run. We commit to
running both before any claim in this paper is treated as out-of-sample
robust. The framework is the contribution; the candidate book is its current
output, not its endpoint.

\bibliography{references}

\appendix

\section{Per-strategy result tables (full leaderboard)}
\label{app:leaderboard}
Table~\ref{tab:leaderboard} (in \S\ref{sec:empirical}) is the full
common-window leaderboard. Per-strategy result cards live in the companion
repository under
\texttt{freqtrade-experiment/hmm-slope-experiment/research/analysis/reports/2026-05-*.md}. Each card contains: trade
log summary, per-layer metric table, MDB triplet, correlation row, and the
pre-registered kill-criteria document referenced.

\section{MDB definition and computation}
\label{app:mdb}

Given strategies $\{s_1, \ldots, s_k\}$ with daily log-return series
$\{r^{(i)}_t\}_{t=1}^T$, define the weighted book return as
$r^{B,w}_t = \sum_i w_i \, r^{(i)}_t$ and its annualised Sharpe as
\begin{equation}
\mathrm{SR}^{B,w} = \sqrt{365} \cdot \frac{\bar{r}^{B,w}}{\sigma^{B,w}}.
\end{equation}

\textbf{Equal weights:} $w_i = 1/k$ trivially.

\textbf{Risk-parity weights} are computed on a trailing 90-day return volatility
$\sigma^{(i)}_t$ at each $t$, with $w^{(i)}_t \propto 1/\sigma^{(i)}_t$ and
renormalised so $\sum_i w^{(i)}_t = 1$. The reported MDB-rp uses the
endpoint weights for clarity; we have verified that re-computing under
rolling weights changes the headline number by less than $0.02$.

\textbf{Mean-variance weights} solve
\begin{align}
\max_w \;& \, w^\top \mu \;-\; \tfrac{1}{2} \lambda \, w^\top \Sigma w \\
\text{s.t.} \;& \, w_i \geq 0,\; \sum_i w_i = 1,
\end{align}
with $\mu, \Sigma$ the sample mean and covariance of daily returns and
$\lambda = 1$. We acknowledge this is unstable at our $N$ and report
MDB-mv as an upper bound only; the robust flag (positive under all three
schemes) is the binding gate.

The common window is the intersection of all strategy backtest date ranges
($\geq 90$ days required). Missing days within a strategy's nominal range
are treated as flat ($r = 0$); we have checked that switching to
NaN-and-pairwise-complete shifts pairwise correlations by less than $0.05$
on average.

\section{Kill-criteria pre-registration template}
\label{app:template}

Every new strategy in the universe receives a dated decision document under
\texttt{freqtrade-experiment/hmm-slope-experiment/research/analysis/reports/YYYY-MM-DD-decision-00N-kill-criteria-<strategy>.md}
\emph{before} its first backtest. The template (compressed):

\begin{enumerate}
  \item Strategy name, family code, universe, timeframe.
  \item Hypothesis (one sentence).
  \item Reference backtest baseline (if reusing an existing measurement) or
        \texttt{N/A} if first-of-family.
  \item \textbf{Hard kills (K1--K4):} hard-MDD ceiling, consecutive
        full-stop trades, rolling-365-day return floor, walk-forward
        quarterly Calmar floor. Each specified as a scalar with rationale
        tied to either the reference baseline or a family-level prior.
  \item \textbf{Continuous shrinkage formula} (Davies--Ravagnani style):
        size multiplier as a product of factor terms, each clipped to
        $[0, 1]$.
  \item \textbf{Re-evaluation triggers:} regime shifts, fee-schedule changes,
        funding-sign flips, cross-family comparisons. Trigger an explicit
        written review within 5 trading days; outcome must be one of
        \{continue at current size, force-shrink, hard-kill\}.
  \item Source-of-truth note: this document wins over any conflicting
        leaderboard cell until superseded by a new dated decision document.
\end{enumerate}

Decision 004~\citep{decision_004} is the canonical worked example. Documents
006, 007, 008 follow the same template for family-specific K1 thresholds
(pairs $K1 = 8\%$, cross-sectional $K1_{\mathrm{xs}} = 12\%$,
funding-MR $K1_{\mathrm{fmr}} = 5\%$).

\end{document}
